\documentclass[UTF8,11pt,a4paper]{ctexart}

\usepackage{amsmath,amssymb,mathtools}
\usepackage{geometry}
\usepackage{enumitem}
\usepackage{array,booktabs,tabularx}
\usepackage{fancyvrb}
\usepackage[hidelinks]{hyperref}

\geometry{top=2.3cm,bottom=2.3cm,left=2.5cm,right=2.5cm}
\setlength{\parindent}{2em}
\setlength{\parskip}{0.25em}
\setlength{\emergencystretch}{2em}
\setlist[itemize]{leftmargin=2em,itemsep=0.45em,topsep=0.35em}
\setlist[enumerate]{leftmargin=2.2em,itemsep=0.45em,topsep=0.35em}

\begin{document}


\section{问题与动机}

对于数组而言，其最重要的性质是随机访问：给定下标 $i$，只需根据首地址和元素宽度计算目标地址，便能在 $\mathrm{O}(1)$ 时间内读取或修改元素。然而，传统连续数组的容量在分配时已经固定。若数组之后需要增长，紧随原数组之后的地址可能已被其他对象占用，算法便不能保证原地延长原有内存区间。因此，可变长数组必须借助动态内存管理。\par

动态内存管理不只意味着申请一个更大的连续数组。其实现至少有两条基本路线：\par

\begin{enumerate}
\item \textbf{连续重分配。} 申请一个更大的连续数组，把旧元素复制过去，再释放旧数组；几何扩缩容属于这一路线。
\item \textbf{分块表示。} 把一个逻辑数组分散到多个动态申请的固定长度数据块中，再用索引块保存各数据块的指针；HAT、Brodnik 结构以及 Tarjan-Zwick 的新构造都采用这一路线。
\end{enumerate}

无论采用哪条路线，外部看到的仍是同一个按下标排列的逻辑数组。区别只在于元素是否物理连续，以及实现如何在更新时间、索引开销、未使用槽位和重组峰值之间取舍。\par

本文研究的抽象数据类型支持 \textnormal{\texttt{Get(i)}}、\textnormal{\texttt{Set(i,x)}}、\textnormal{\texttt{Grow(x)}} 和 \textnormal{\texttt{Shrink()}}。这里的“末端”或论文所说的“远端”，是指当前最大下标一侧：若数组原来包含 \textnormal{\texttt{a0,a1,...,a(N-1)}}，\textnormal{\texttt{Grow(x)}} 产生的新元素是 $aN=x$，\textnormal{\texttt{Shrink()}} 删除的是 \textnormal{\texttt{a(N-1)}}。因此，末端更新不会改变任何保留元素的下标。\par

这一操作边界与任意位置插入、删除的动态序列有本质区别。例如，在\par

\begin{Verbatim}[fontsize=\scriptsize]
[a0, a1, a2, a3]
\end{Verbatim}

的位置 1 插入 \textnormal{\texttt{x}} 后，序列变成\par

\begin{Verbatim}[fontsize=\scriptsize]
[a0, x, a1, a2, a3].
\end{Verbatim}

原来的 \textnormal{\texttt{a1,a2,a3}} 都获得了新下标，数据结构不仅要容纳一个新元素，还要维护后续元素变化后的秩。本文的末端 \textnormal{\texttt{Grow/Shrink}} 不产生这种全局秩变化，所以不能只比较大 O 数值就把两类问题视为相同。论文的单端结果也可以组合成双端可变长数组：分别用两个单端结构保存中心点左右两侧，其中一侧按相反顺序编号。双端更新仍只作用于两个物理末端，也不等于任意位置插删。\par

最常见的连续重分配方案是几何扩缩容。它保持最坏 $\mathrm{O}(1)$ 的随机访问和摊还 $\mathrm{O}(1)$ 的末端更新，却可能长期保留 $\Theta(N)$ 空槽，并在重建时让新旧两个线性大小的数组短暂共存。HAT 和 Brodnik 等分块方案把冗余降低到 $\mathrm{O}(\sqrt{N})$，但必须额外维护指针和定位规则。\par

不过，如果仅写“某个方案使用 $\mathrm{O}(N)$ 空间”，则会掩盖两个不同问题：数组在一次操作完成后需要长期保留多少空间，以及一次扩缩容正在执行时会短暂达到多高的空间峰值。Tarjan 和 Zwick 的关键出发点，就是分别度量这两个量。\par

\section{模型与双空间指标}

\subsection{内存块模型}

根据论文原文所述，其将具体实现视为一组动态申请的固定长度内存块（blocks）。数据块保存元素，索引块保存指针，辅助块保存长度或计数。若当前数组有 $N$ 个机器字大小的元素，则数据本身至少需要 $N$ 个机器字；除一个主索引块外，每个可访问内存块还需要被某个指针指向。这里的“重新分配”和“重建”含义不同：重新分配（\textnormal{\texttt{Reallocate}}）是内存管理接口对单个连续块进行扩展、搬移或替换；重建（rebuild）是数据结构重新组织多个数据块、索引块和元素的过程。一次重建可能调用多次申请、释放或重新分配。\par

\subsection{稳定空间与扩缩容临时空间}

对两个非递减函数 $s(N)$ 和 $t(N)$，若一个实现满足：\par

\begin{Verbatim}[fontsize=\scriptsize]
稳定保存大小为 N 的数组：最多 N+s(N) 空间；
执行 Grow 或 Shrink 期间：最多 N+t(N) 空间，
\end{Verbatim}

则称其为 \textnormal{\texttt{(s(N),t(N))}}-实现。本文把 $s(N)$ 统一称为稳定冗余，把 $N+s(N)$ 称为稳定存储空间；把 $t(N)$ 称为扩缩容临时冗余，把 $N+t(N)$ 称为扩缩容临时空间。“稳定”表示一次操作完成后的表示，不表示这些内存永久不释放；“临时”强调 \textnormal{\texttt{Grow/Shrink}} 的执行过程，而不是另一种独立的数据结构。\par

双指标的意义在于，应用可能同时保存许多可变长数组，却只在某一时刻调整其中少数数组。此时大多数数组可以保持紧凑表示，共享的工作内存只为正在重组的结构提供短暂空间。不过，这只是理论动机；实际收益仍受缓存、内存分配器、碎片和并发重组影响。\par

\subsection{全文术语约定}

为避免把不同论文的局部名称误认为不同概念，本文采用以下主要术语：\par

\begin{center}
\small
\begin{tabularx}{\textwidth}{@{}>{\raggedright\arraybackslash}X>{\raggedright\arraybackslash}X>{\raggedright\arraybackslash}X@{}}
\toprule
\textbf{原文或局部名称} & \textbf{本文主要名称} & \textbf{使用规则} \\
\midrule
\textnormal{\texttt{Top}} & 索引块 & 介绍 Sitarski 原文时写作“Top（索引块）” \\
\textnormal{\texttt{Leaf}} & 数据块 & 介绍 Sitarski 原文时写作“Leaf（数据块）” \\
\textnormal{\texttt{virtual superblock}} & 虚拟超级块 & 仅表示 Brodnik 对同尺度数据块的概念分组，不是实际申请的内存块 \\
\textnormal{\texttt{Reallocate}} & 重新分配 & 指内存接口对一个块的操作 \\
\textnormal{\texttt{rebuild}} & 重建 & 指数据结构层面的重新组织 \\
\textnormal{\texttt{access}} & 访问 & 包括 \textnormal{\texttt{Get}} 与 \textnormal{\texttt{Set}} 所需的下标定位和读写 \\
\textnormal{\texttt{amortized time}} & 摊还时间 & 针对任意操作序列的总成本，不写成“平均时间” \\
\textnormal{\texttt{standard implementation}} & 标准实现 & 仅在 Sections 7-9 的作者定义范围内使用 \\
\textnormal{\texttt{growth game}} & 增长博弈 & 指时间下界所用的抽象博弈 \\
\bottomrule
\end{tabularx}
\end{center}

有了操作模型和双空间指标，下面可以把旧方法放到同一坐标系中比较：精确容量方案追求最小稳定冗余却频繁复制，几何扩缩容以线性冗余换取简单和快速，HAT 与 Brodnik 则把任意时刻的冗余降到平方根级。\par

\section{三类旧方法}

\subsection{精确容量与几何扩容}

若每次 \textnormal{\texttt{Grow}} 或 \textnormal{\texttt{Shrink}} 都申请恰好匹配新长度的连续数组，稳定冗余只有 $\mathrm{O}(1)$，但每次更新都要复制 $\Theta(N)$ 个元素，重建时总空间接近 $2N$。它展示了极小稳定空间的一端。\par

几何扩容位于另一端。设固定扩容因子为\par

\[
  g = 1+\alpha,  \alpha>0.
\]

当容量为 $N$ 的连续数组刚好装满时，结构申请一个容量约为\par

\[
  (1+\alpha)N
\]

的新数组，复制原来的 $N$ 个元素，然后释放旧数组。忽略取整和触发扩容的单个新元素，重建刚结束时大约新增 \textnormal{\texttt{alpha N}} 个空槽，下一次装满前也会发生大约 \textnormal{\texttt{alpha N}} 次 \textnormal{\texttt{Grow}}。一次重建复制 $N$ 个元素，因此分摊到这些更新上的复制成本为\par

\[
  N/(\alpha N) = 1/\alpha
\]

也就是 $\Theta(1/\alpha)$。若把每次 \textnormal{\texttt{Grow}} 自身写入新元素的常数成本也计入，论文给出的表达式与 $(1+\alpha)/\alpha$ 同阶。因而 $\alpha$ 变小会产生一组方向相反的影响：\par

\begin{Verbatim}[fontsize=\scriptsize]
alpha 变小
-> 每次扩容预留的空槽减少
-> 两次扩容之间可容纳的 Grow 数量减少
-> 扩容更频繁
-> 每次 Grow 的摊还复制成本增大
\end{Verbatim}

两个数字例子可以显示这一差别：\par

\begin{center}
\small
\begin{tabularx}{\textwidth}{@{}>{\raggedright\arraybackslash}X>{\raggedright\arraybackslash}X>{\raggedright\arraybackslash}X>{\raggedright\arraybackslash}X>{\raggedright\arraybackslash}X@{}}
\toprule
\textbf{增长因子 $g$} & \textbf{$\alpha$} & \textbf{刚扩容后的新增空槽} & \textbf{约需多少次 \textnormal{\texttt{Grow}} 再次装满} & \textbf{分摊复制量级} \\
\midrule
2 & 1 & $N$ & $N$ & 约 1 次复制/操作 \\
1.25 & 0.25 & $0.25N$ & $0.25N$ & 约 4 次复制/操作 \\
\bottomrule
\end{tabularx}
\end{center}

这里“约 1 次”或“约 4 次”指元素复制次数的摊还量级，不表示一次操作的实际墙钟时间，也忽略了取整、分配器和元素构造成本。增长因子为 2 时，空闲槽数最多可接近当前元素数，也就是已分配容量可能接近一半未使用；增长因子为 1.25 时，刚扩容后的空闲槽约占新容量的 $0.25/1.25=20\%$，但重建明显更频繁。\par

缩容不能简单地在使用率刚低于 $1/g$ 时立即进行。否则数组在边界附近执行一次 \textnormal{\texttt{Shrink}} 后缩容，紧接着一次 \textnormal{\texttt{Grow}} 又可能重新扩容，两个昂贵重建无法由足够多的普通操作分摊。标准做法是使用滞后阈值：若当前容量为 $C$，可以等到元素数降至约\par

\[
  C/g^2
\]

时，才把容量缩小到约 $C/g$。缩容后装载率约为 $1/g$，离再次扩容或缩容都保留了一段缓冲区间。这个滞后区间保证相邻重建之间发生线性数量的更新，从而维持 $\mathrm{O}(1)$ 摊还成本。\par

对任意固定 \textnormal{\texttt{alpha>0}}，几何方案在稳定状态中保留的空槽最多为 $\Theta(\alpha N)=\Theta(N)$，而重建时旧数组和容量约为 $(1+\alpha)N$ 的新数组同时存在，临时冗余也是 $\Theta(N)$。采用后台复制可以把一次大规模重建分散到后续操作，从而获得最坏情况常数更新时间，但这不会自动消除新旧表示共存所需的线性空间。由此得到几何扩缩容的完整取舍：它用线性冗余换取简单地址计算和常数更新界，而减小增长因子只能改善常数比例，不能把冗余降为 \textnormal{\texttt{o(N)}}。\par

线性冗余并不是常数更新时间的必要代价。Sitarski 的 HAT 改用等长数据块，只在较少的尺度变化时重建，从而把下面要讨论的冗余降到平方根级。\par

\subsection{Sitarski 的 HAT}

即便名字中包含"哈希"二字，Sitarski 的哈希数组树（hashed array tree，HAT）并不使用哈希，结构也不是一棵通常意义上的树，而是一个索引块指向若干等长数据块的两级表示。对当前长度 $N$，结构维护一个通常取为 2 的幂的参数 $B$，并保持\par

\[
  \sqrt{N} \leq B < 4\sqrt{N}.
\]

HAT 的结构不变量如下：\par

\begin{itemize}
\item 一个长度为 $B$ 的索引块 \textnormal{\texttt{I}}，其中每个已使用位置保存一个数据块指针；
\item $\lceil N/B\rceil$ 或 $\lceil N/B\rceil+1$ 个长度均为 $B$ 的数据块；
\item 前面的数据块全部填满，只有最后一个非空数据块可以部分填充；
\item 结构至多保留一个完全空的数据块，以避免一次 \textnormal{\texttt{Shrink}} 刚释放末块、下一次 \textnormal{\texttt{Grow}} 又立即重新申请；
\item 所有逻辑元素仍按下标顺序依次排列在这些数据块中。
\end{itemize}

因此，对从 0 开始的下标 $i$，元素位置可直接计算为\par

\[
\begin{aligned}
  \operatorname{block}(i)=\left\lfloor\frac{i}{B}\right\rfloor \\
  \operatorname{offset}(i)=i\bmod B.
\end{aligned}
\]

访问过程先读取 \textnormal{\texttt{I[block(i)]}} 得到数据块指针，再读取该块的 \textnormal{\texttt{offset(i)}} 位置，只进行常数次内存访问和整数运算。若 $B$ 是 2 的幂，除法和取模还可分别由右移和位掩码实现，所以 \textnormal{\texttt{Get}} 与 \textnormal{\texttt{Set}} 都是最坏 $\mathrm{O}(1)$。\par

\textnormal{\texttt{Grow}} 在当前尺度尚有容量时分为三种情况：\par

\begin{enumerate}
\item 最后一个数据块未满，直接把新元素写入其下一个空位；
\item 最后一个数据块已满，但结构已经保留一个空数据块，把新元素写入该空块首位；
\item 最后一个数据块已满且没有预留空块，申请一个新的长度 $B$ 数据块，在索引块中登记其指针，再写入新元素。
\end{enumerate}

当 $N=B^2$，长度为 $B$ 的索引块已经指向 $B$ 个满数据块，当前尺度达到容量上限。结构把参数加倍为 $2B$ 并重建。\textnormal{\texttt{Shrink}} 的普通情况与 \textnormal{\texttt{Grow}} 对称：删除最后一个元素；末端块变空后可以暂不释放，只有在末端出现两个空块时才释放其中一个。当数组缩小到 $N=B^2/16$，等价于当前 $B=4sqrt(N)$，结构把 $B$ 减半并重建。扩张和收缩使用不同阈值，避免操作序列在一个边界附近来回触发重建。\par

稳定状态的冗余可以逐项计算：\par

\begin{Verbatim}[fontsize=\scriptsize]
索引块                         B
最后一个部分填充数据块的空位   < B
至多一个预留空数据块            B
长度、计数等辅助信息            O(1)
\end{Verbatim}

因此总空间至多为 $N+3B+\mathrm{O}(1)=N+\mathrm{O}(\sqrt{N})$。改变 $B$ 时不必同时申请另一套包含 $N$ 个槽的完整结构；新数据块可以逐个申请、填充，旧数据块也可逐个释放，所以重建只需 $\mathrm{O}(\sqrt{N})$ 级临时冗余。另一方面，加倍或减半 $B$ 会使可表示容量改变常数倍，两次同方向重建之间至少发生 $\Omega(N)$ 次 \textnormal{\texttt{Grow}} 或 \textnormal{\texttt{Shrink}}，故一次 $\mathrm{O}(N)$ 重建可摊到这些操作上，每次更新的摊还代价仍为 $\mathrm{O}(1)$。\par

下面是一个自行绘制的小例子。令 $B=4,N=13$，索引块有 4 个位置，前三个数据块填满，第四个数据块只含元素 \textnormal{\texttt{a12}}：\par

\begin{Verbatim}[fontsize=\scriptsize]
index block I (length 4)
+------+------+------+------+
|  D0  |  D1  |  D2  |  D3  |
+--|---+--|---+--|---+--|---+
   |      |      |      |
   v      v      v      v
 D0       D1       D2       D3
+---+---+---+---+  +---+---+---+---+  +---+---+---+---+  +---+---+---+---+
|a0 |a1 |a2 |a3 |  |a4 |a5 |a6 |a7 |  |a8 |a9 |a10|a11|  |a12|   |   |   |
+---+---+---+---+  +---+---+---+---+  +---+---+---+---+  +---+---+---+---+
\end{Verbatim}

例如 $i=10$ 时，$\lfloor 10/4\rfloor=2$ 且 $10\bmod 4=2$，所以 \textnormal{\texttt{a10}} 位于 \textnormal{\texttt{D2}} 的偏移 2。这个例子也直接显示：末块浪费 3 个槽，但浪费始终小于 $B$。\par

HAT 说明可变长数组并不必须浪费常数比例空间：分块和一级索引能够把冗余降到 $\mathrm{O}(\sqrt{N})$，同时保留常数访问和常数摊还更新。它仍会在参数 $B$ 改变时重建数据布局；Brodnik 等人的结构进一步改变块尺度组织，使数据元素不必周期性整体搬移。\par

\subsection{Brodnik 等人的结构}

在Sitarski 的哈希数组树之外，Brodnik、Carlsson、Demaine、Munro 与 Sedgewick 提出了另一种 $N+\mathrm{O}(\sqrt{N})$ 的可变长数组。论文给出两个数据布局变体：第一种使用严格递增的数据块；第二种把数据块大小限制为 2 的幂，以便避免访问定位中的平方根计算。\par

第一种变体把元素按顺序放入容量为\par

\begin{Verbatim}[fontsize=\scriptsize]
1, 2, 3, ..., k
\end{Verbatim}

的数据块。前 $k$ 个块的总容量为三角数\par

\[
  1+2+\ldots+k = k(k+1)/2.
\]

因此，足以保存 $N$ 个元素的最小块数为\par

\[
  k=\left\lceil\frac{\sqrt{8N+1}-1}{2}\right\rceil=\Theta(\sqrt{N}).
\]

除最后一个块外，其余块都填满；最后一个块可以部分填充。结构还可以保留一个曾经含有元素、后来因 \textnormal{\texttt{Shrink}} 变空的下一数据块。一个大小为 $\Theta(\sqrt{N})$ 的索引块保存这些数据块的指针，而索引块自身使用最朴素的动态数组方法扩展。数据块、尾块空位、数据块指针以及索引块的未用位置都只有 $\mathrm{O}(\sqrt{N})$ 量级，所以稳定总空间为 $N+\mathrm{O}(\sqrt{N})$。\par

该布局的 \textnormal{\texttt{Grow}} 流程是：\par

\begin{enumerate}
\item 若最后一个非空数据块尚未填满，直接写入其下一个位置；
\item 若该块已满但已经存在一个预留空块，把新元素写入空块首位；
\item 否则申请容量为下一整数的新数据块，并把指针加入索引块；
\item 若索引块已满，则申请更大的索引块并复制全部数据块指针；若要求最坏 $\mathrm{O}(1)$ 更新时间，这次指针复制需在后续操作中后台完成。
\end{enumerate}

\textnormal{\texttt{Shrink}} 删除最后一个元素。若末端数据块变空，可先把它作为预留空块保存；当继续收缩而不再需要它时再释放，并从索引块删除相应指针。这里移动的主要是索引指针，已经装入数据块的元素不会因为全局尺度变化而被周期性整体搬迁。这是它与 HAT 的重要区别。\par

第一种变体的访问定位来自反解三角数。为避免 0/1 下标混淆，令 $p=i+1$ 表示元素的 1-based 位置，令 $j$ 表示 1-based 数据块编号，则 $j$ 是满足 $j(j+1)/2\geq p$ 的最小整数，即\par

\[
  j=\left\lceil\frac{\sqrt{8p+1}-1}{2}\right\rceil.
\]

该元素在第 $j$ 个数据块中的 0-based 偏移为\par

\[
  offset = p - j(j-1)/2 - 1.
\]

例如 $N=8$ 时，块容量依次为 \textnormal{\texttt{1,2,3,4}}，第四块只使用前两个位置。访问 $i=6$，即 $p=7$，可得 $j=4$、$offset=0$，所以该元素位于第四块首位。\par

\begin{Verbatim}[fontsize=\scriptsize]
index block
+------+------+------+------+
|  D1  |  D2  |  D3  |  D4  |
+--|---+--|---+--|---+--|---+
   v      v      v      v
+----+  +----+----+  +----+----+----+  +----+----+----+----+
| a0 |  | a1 | a2 |  | a3 | a4 | a5 |  | a6 | a7 |    |    |
+----+  +----+----+  +----+----+----+  +----+----+----+----+
 size 1    size 2         size 3              size 4
\end{Verbatim}

直接计算上式需要整数平方根。为避免依赖该操作，第二种变体只使用大小为 2 的幂的数据块，并把它们组织成虚拟超级块（virtual superblocks）：第 $q$ 个虚拟超级块由 $2^{\lfloor q/2\rfloor}$ 个、每个容量为 $2^{\lceil q/2\rceil}$ 的数据块组成。虚拟超级块只是对相同尺度数据块的逻辑分组，并没有对应的整块内存申请。由元素下标的最高有效位和若干移位即可确定虚拟超级块、其中的数据块以及块内偏移。正文不需要复现其全部位级公式，但必须保留这一逻辑：第一种变体通过反三角数定位，第二种变体通过幂次分组把定位转化为 word-RAM 字操作；在相应机器运算假设下，两者都支持最坏 $\mathrm{O}(1)$ 访问。\par

HAT 与 Brodnik 结构的差异可以概括如下：\par

\begin{center}
\small
\begin{tabularx}{\textwidth}{@{}>{\raggedright\arraybackslash}X>{\raggedright\arraybackslash}X>{\raggedright\arraybackslash}X@{}}
\toprule
\textbf{方面} & \textbf{HAT} & \textbf{Brodnik 等人的结构} \\
\midrule
数据块大小 & 全部为当前参数 $B$ & 递增大小，或按 2 的幂组织 \\
下标定位 & 商和余数 & 反三角数，或最高有效位与移位 \\
数据元素重建 & $B$ 改变时需要重建数据布局 & 不需要周期性整体搬移数据元素 \\
索引处理 & 改变 $B$ 时一并重建 & 索引块可独立扩展并后台复制 \\
空间 & $N+\mathrm{O}(\sqrt{N})$ & $N+\mathrm{O}(\sqrt{N})$ \\
更新界 & $\mathrm{O}(1)$ 摊还；可进一步去摊还化 & 在分配假设和后台索引复制下最坏 $\mathrm{O}(1)$ \\
\bottomrule
\end{tabularx}
\end{center}

两种方案从不同结构出发，却共同达到 $N+\mathrm{O}(\sqrt{N})$。这使问题自然转向下界：平方根是否只是两个设计恰好得到的结果，还是单一峰值空间指标下不可突破的屏障？Brodnik 等人的论证回答了后一个问题。\par

\section{为什么旧下界仍允许新结果}

有别于新下界将稳定状态和重建瞬间严格区分，旧下界把稳定状态和重建瞬间统一计入“任意时刻空间”，从而使得旧下界需在所有时刻严格保持 $N+\mathrm{O}(\sqrt{N})$。Tarjan 和 Zwick 观察到，$N+\Omega(\sqrt{N})$ 只要求某个历史时刻出现该冗余，而这个时刻可能恰好是新块刚被申请、旧表示尚未释放的扩缩容过程。因而，只要接受扩缩容期间更高的临时峰值，一个结构仍可能在操作结束后只保留远少于 $\sqrt{N}$ 的稳定冗余。而这一点可以被推广成稳定冗余与临时冗余的乘积下界。\par

\section{空间下界}

为便于理解，下面两个证明使用同一组内存模型假设。数组元素是任意 word-size 值，因而不能默认进一步压缩；每个数据块都必须由一个指针标识，一个指针占 $\Theta(1)$ words。索引块与长度、计数等辅助信息也计入空间，而且只会消耗更多冗余，不会使可用数据块数量增加。最后，新块刚申请时为空；在元素复制完成以前，旧元素仍必须存在于原数据块中。Theorem 4.2 还使用 Definition 2.1 的约定：$s(N)$ 和 $t(N)$ 都是非递减函数。\par

\subsection{\texorpdfstring{Theorem 4.1：$\sqrt{N}$ 下界}{Theorem 4.1：sqrt(N) 下界}}

\textbf{定理。} 即使操作序列只包含 \textnormal{\texttt{Grow}} 和访问，任何可变长数组实现也必定在某些时刻使用 $N+\Omega(\sqrt{N})$ 空间，其中这里的 $N$ 指该时刻数组的长度。\par

\textbf{证明。} 从空数组开始执行一段共 $N$ 次的 \textnormal{\texttt{Grow}} 操作。执行完毕时，$N$ 个元素分布在 $k$ 个连续数据块中。根据 $k$ 的大小分两种情况。\par

若 $k\geq \sqrt{N}$，保存数据本身至少需要 $N$ words，而标识这 $k$ 个数据块至少需要 $k$ 个指针。因此最终时刻的空间至少为\par

\[
  N+k \geq  N+\sqrt{N}.
\]

若 $k<\sqrt{N}$，由平均值原理，至少存在一个容量为 $\ell$ 的数据块满足\par

\[
  \ell \geq  N/k > \sqrt{N}.
\]

设这个大块在历史上数组长度为 $N'\leq N$ 时被申请。新块刚申请时还没有替代旧表示，因此当时已有的 $N'$ 个元素仍在其他数据块中，新块的 $\ell$ 个槽必须与它们短暂共存。又因为 $N'\leq N$，所以\par

\[
  \ell > \sqrt{N} \geq  \sqrt{N'}.
\]

故该历史时刻的总空间至少为\par

\[
  N'+\ell > N'+\sqrt{N'}.
\]

在定理陈述中把这个发生峰值时的数组长度 $N'$ 重新记为 $N$，便得到“某些时刻使用 $N+\Omega(\sqrt{N})$ 空间”的结论。证毕。\par

证明中的两个 case 可以概括为：块多会增加指针开销；块少则迫使某个块很大，而这个大块被申请时必须与旧元素共存。需要注意，定理没有声称每个稳定状态都有 $\sqrt{N}$ 冗余；第二种情况的峰值可能发生在早于最终第 $N$ 次 \textnormal{\texttt{Grow}} 的历史时刻。该下界也不依赖 \textnormal{\texttt{Grow}} 究竟花费多少时间。\par

\subsection{Theorem 4.2：乘积下界}

\textbf{定理。} 任何 \textnormal{\texttt{(s(N),t(N))}}-implementation 都必须满足\par

\[
  s(N)t(N) \geq  N
\]

即使只支持 \textnormal{\texttt{Grow}} 和访问。\par

\textbf{证明。} 在第 $N$ 次 \textnormal{\texttt{Grow}} 完成后的稳定状态，额外空间至多为 $s(N)$。由于每个数据块至少需要一个指针，并且每个指针占一个常数数量的 words，数据块数至多为 $s(N)$（若精确考虑指针表示常数，则结论只相差常数因子）。索引块的空槽和其他辅助信息同样占用 $s(N)$，只会让实际可保存的指针数更少，不会破坏这一上界。\par

$N$ 个元素分布在至多 $s(N)$ 个数据块中。由平均值原理，至少有一个数据块 $B$ 的容量满足\par

\begin{Verbatim}[fontsize=\scriptsize]
|B| >= N/s(N).
\end{Verbatim}

设 $B$ 在历史上数组长度为 $N'\leq N$ 时被申请。刚申请时 $B$ 为空，复制完成前原有元素仍在其他数据块中，所以这一时刻的临时额外空间至少为 \textnormal{\texttt{|B|}}，即\par

\begin{Verbatim}[fontsize=\scriptsize]
t(N') >= |B| >= N/s(N).
\end{Verbatim}

现在才使用 \textnormal{\texttt{t}} 的非递减性。由于 $N'\leq N$，\par

\[
  t(N) \geq  t(N') \geq  N/s(N)
\]

两边乘以 $s(N)$，得到\par

\[
  s(N)t(N) \geq  N.
\]

证毕。\par

这个证明中，块数上界来自指针空间，最大块下界来自平均值原理，而单调性只用于把历史时刻的 \textnormal{\texttt{t(N')}} 转换成最终参数 $t(N)$。它证明的是稳定冗余与扩缩容临时冗余之间的空间权衡，并没有推出任何 $\Omega(r)$ 更新时间下界。\par

\textbf{Corollary 4.3。} 若\par

\[
  s(N)=\mathrm{O}(N^{1/r})
\]

则 Theorem 4.2 给出\par

\[
\begin{aligned}
  t(N) = \Omega(N/s(N)) \\
  = \Omega(N^{1-1/r}).
\end{aligned}
\]

两个最重要的具体代入为：\par

\begin{center}
\small
\begin{tabularx}{\textwidth}{@{}>{\raggedright\arraybackslash}X>{\raggedright\arraybackslash}X>{\raggedright\arraybackslash}X@{}}
\toprule
\textbf{$r$} & \textbf{稳定冗余 $s(N)$} & \textbf{必需的临时冗余下界 $t(N)$} \\
\midrule
2 & $\mathrm{O}(N^{1/2})$ & $\Omega(N^{1/2})$ \\
3 & $\mathrm{O}(N^{1/3})$ & $\Omega(N^{2/3})$ \\
\bottomrule
\end{tabularx}
\end{center}

$r=2$ 对应 HAT 和 Brodnik 等旧结构所处的平衡点；$r=3$ 则预告 Section 5 的新构造。如果希望稳定状态只有立方根级冗余，某些 \textnormal{\texttt{Grow}} 中就不可避免地出现三分之二次幂级临时冗余。\par

\section{与论文后续内容的连接}

Section 5 的 $r=3$ 结构正好达到 $N+\mathrm{O}(N^{1/3})$ 稳定空间和 $N+\mathrm{O}(N^{2/3})$ 临时空间；Section 6 先得到 $N+\mathrm{O}(rN^{1/r})$ 稳定空间，Section 7 再对 $r(N) \leq  (1/2) \log N / \log \log N$ 给出 $N+\mathrm{O}(N^{1/r})$ 稳定空间。因此 Theorem 4.2 不只是一个独立的限制，它预先给出了后续构造应该追求的指数。对固定整数 $r$，上述增长条件自然成立。\par

然而，空间下界并不说明 \textnormal{\texttt{Grow/Shrink}} 的摊还时间必须为 $\Omega(r)$。该时间下界需要 Sections 8-9 的增长博弈，并且严格适用于作者定义的标准实现。最终 review 应把“空间最优”和“时间最优”分开证明，再在结论处合并，不能只引用 Theorem 4.2 就声称整篇论文已经在所有维度最优。\par

\section{相关工作与发表后发展}

前文为了验证复杂度，已经技术性地介绍了几何扩缩容、HAT 和 Brodnik 结构。本节不重复其完整操作流程，而是回答三个不同问题：这些结构为何被提出，它们在更大的动态数据结构研究中处于什么位置，以及 Tarjan-Zwick 发表后出现了哪些实现和经验工作。\par

\subsection{可变长数组的直接前驱}

标准意义下的几何扩缩容确立了最常见的基线：连续存储带来最简单的常数访问；固定增长因子带来常数摊还更新；代价是稳定状态与重建过程都可能出现 $\Theta(N)$ 冗余。减小增长因子只能改变常数比例，不能得到 $N+o(N)$ 空间。\par

Sitarski 1996 年提出 HAT 的直接动机来自数据库程序：查询结果长度无法预知，商业可变长数组类在扩展时执行大量复制。其原文不仅给出 Top（索引块）和 Leaves（数据块）的两级结构，还分析了位级寻址、内存碎片和累计复制量。对于 $N=4^m$ 个顺序加入的元素，重建复制量形成 $1+4+\ldots+N=(4N-1)/3$，因此总复制为 $\mathrm{O}(N)$。原文还估计最坏冗余约为 $2sqrt(N)$，并在 1990 年代硬件上比较 HAT 与普通 C++ 数组。由此可见，HAT 从一开始就同时关注渐近空间、复制次数和工程可用性；但其旧实验只能作为历史证据，不能直接预测现代机器性能。\par

Brodnik 等人 1999 年从随机队列（randomized queue）出发，把单端可变长数组放入更完整的动态分配模型中。原文把内存块头部（headers）计入空间，定义 $Allocate/Deallocate/Reallocate$，并用最大块大小 \textnormal{\texttt{f(N)}} 与块数量 \textnormal{\texttt{g(N)}} 的关系\par

\[
  f(N)g(N) \geq  N
\]

得到 $\Omega(\sqrt{N})$ 峰值冗余下界。其最终结构使用虚拟超级块和位运算定位，并给出可执行的 \textnormal{\texttt{Grow}}、\textnormal{\texttt{Shrink}} 与 \textnormal{\texttt{Locate}} 算法；通过后台复制索引指针得到最坏 $\mathrm{O}(1)$ 操作，还给出适配伙伴系统（buddy system）的版本。该结构进一步导出栈、队列、随机队列、优先队列和双端队列等结果。\par

两项旧工作共同优化的是“任意时刻总空间”，其自然平衡点为 $N+\Theta(\sqrt{N})$。Tarjan-Zwick 并未否定这一旧下界，而是改变了度量：把稳定存储空间和扩缩容临时空间分开，再研究两者与更新时间的完整权衡。\par

\subsection{简洁动态数据结构背景}

在简洁动态数据结构层面，可变长数组的低冗余表示属于简洁数据结构（succinct data structures）的广泛背景。Raman、Raman 与 Rao 研究动态简洁结构，Raman 与 Rao 进一步研究动态字典和树，Munro 与 Rao 的综述则系统整理了相关表示方法。这个方向通常要求数据结构占用接近信息论最低编码长度的比特数，并在低阶附加空间内支持查询和更新。\par

Tarjan-Zwick 与该方向共享“主数据加低阶冗余”的目标，但空间基线不同：本文把 $N$ 个任意机器字元素本身所需的 $N$ 个机器字作为不可压缩主数据，再优化其上的指针、空槽和重组空间。因此，正文可以把它放入简洁结构研究背景，却不应把 $N+o(N)$ 个机器字与经典的比特级简洁编码当作完全相同的结论。\par

\subsection{更强的一般动态序列问题}

本文只允许在最大下标一端执行 \textnormal{\texttt{Grow/Shrink}}。Dietz 的 list indexing 和 Fredman-Saks 的动态数据结构下界研究的是更强的秩维护问题：若在任意位置插入或删除，后续元素的逻辑下标都会改变。Tarjan-Zwick 引言据此指出，所有操作同时支持时会出现 $\Theta(\log N/\log \log N)$ 量级的界，这与本文的末端常数访问问题不是同一操作模型。\par

Goodrich 与 Kloss 的分层向量（tiered vectors）用常数层分级结构支持一般位置更新；在固定层数下可以保持常数访问，并以 $N$ 的分数次幂时间完成插入和删除。Bille 等人 2017 年进一步实现多层分层向量，支持 $\mathrm{O}(1)$ 访问、$\mathrm{O}(N^epsilon)$ 插删和 \textnormal{\texttt{o(N)}} 冗余，并给出大规模 C++ 实验。这些工作说明“多层分块”也能服务于一般动态序列，但其更新时间不能与只在末端更新的 $\mathrm{O}(r)$ 摊还界直接比较。\par

\subsection{实践研究与发表后发展}

Joannou 与 Raman 对可扩展数组（extensible arrays）进行了经验比较，Katajainen 则研究如何在实践中实现具有最坏情况保证的动态数组。这些研究提醒我们：渐近冗余、访问指令数量、缓存局部性、分配器行为和实现常数是不同评价维度。Sitarski 的早期基准和 Brodnik 的伙伴系统版本也属于这条实践脉络，而不是 Tarjan-Zwick 新下界的一部分。\par

Tarjan-Zwick 的初步版本发表于 SOSA 2023，最终版本发表于 2024 年 \emph{SIAM Journal on Computing}。最终期刊版致谢 Christian Kjær 和 Victor Brevig，并说明他们的实现促使作者简化 Section 6.3 的常数时间访问方法。两人的 DTU 学位论文实现并测试了 Tarjan-Zwick 的结构，补充了原文没有完全展开的重建细节；其实验认为，理论 $\mathrm{O}(1)$ 定位在实际可用的较小 $r$ 上可能慢于简单的 $\mathrm{O}(r)$ 循环定位，并比较了稳定冗余、临时峰值、访问和更新成本。\par

根据尽两年的学术研究，2025 年出现了独立 Rust 实现 \textnormal{\texttt{tzarrays}}，说明该结构已经有论文作者之外的工程实现。不过，学位论文和代码仓库不等同于新的同行评审理论结果。截至 2026-07-16 的本轮公开检索，确认的发表后发展主要是最终期刊版本、实现与经验评估；尚未确认有同行评审理论工作进一步改进其渐近权衡。这是带检索范围和日期的结论，不能改写成“此后绝对没有研究”。\par

\subsection{本文在研究脉络中的位置}

相关工作可以沿三条轴概括。几何扩缩容、HAT 和 Brodnik 逐步降低可变长数组的空间浪费；简洁数据结构提供“主数据加低阶冗余”的一般研究背景；分层向量等结构处理更强的任意位置动态序列。Tarjan-Zwick 的独特位置，是在末端更新模型中首次系统分离稳定存储空间与扩缩容临时空间，并同时匹配空间和更新时间下界。\par

因此，A 部分最终应停在一个明确问题上：当 $r=3$ 时，Corollary 4.3 已经说明 $\mathrm{O}(N^{1/3})$ 稳定冗余必然伴随 $\Omega(N^{2/3})$ 临时冗余；接下来需要由 B 说明，Section 5 如何用大块和小块（large/small blocks）构造一个恰好达到这两个指数、同时保持最坏 $\mathrm{O}(1)$ 访问和摊还 $\mathrm{O}(1)$ 更新的结构。\par


\end{document}
